\documentclass[12pt,a4paper]{article}

%% ─── PAQUETES ────────────────────────────────────────────────────────────────
\usepackage[utf8]{inputenc}
\usepackage[T1]{fontenc}
\usepackage[spanish]{babel}
\usepackage{lmodern}
\usepackage{geometry}
\geometry{margin=2.5cm}
\usepackage{graphicx}
\usepackage{amsmath}
\usepackage{booktabs}
\usepackage{longtable}
\usepackage{array}
\usepackage{xcolor}
\usepackage{hyperref}
\usepackage{fancyhdr}
\usepackage{titlesec}
\usepackage{enumitem}
\usepackage{float}
\usepackage{caption}
\usepackage{subcaption}
\usepackage{listings}
\usepackage{tikz}
\usetikzlibrary{shapes.geometric, arrows.meta, positioning, fit, backgrounds}
\usepackage{mdframed}
\usepackage{setspace}
\usepackage{parskip}

%% ─── COLORES ─────────────────────────────────────────────────────────────────
\definecolor{teal}{RGB}{20,184,166}
\definecolor{emerald}{RGB}{16,185,129}
\definecolor{slate}{RGB}{71,85,105}
\definecolor{lightgray}{RGB}{241,245,249}
\definecolor{codebg}{RGB}{30,41,59}
\definecolor{codetext}{RGB}{226,232,240}

%% ─── HYPERREF ────────────────────────────────────────────────────────────────
\hypersetup{
  colorlinks=true,
  linkcolor=teal,
  urlcolor=teal,
  citecolor=emerald,
  pdftitle={ChatDoc AI -- Agente de Consulta Documental},
  pdfauthor={},
}

%% ─── CABECERA Y PIE ──────────────────────────────────────────────────────────
\pagestyle{fancy}
\fancyhf{}
\fancyhead[L]{\small\textcolor{slate}{ChatDoc AI}}
\fancyhead[R]{\small\textcolor{slate}{Informe Técnico}}
\fancyfoot[C]{\small\thepage}
\renewcommand{\headrulewidth}{0.4pt}

%% ─── TÍTULOS ─────────────────────────────────────────────────────────────────
\titleformat{\section}{\Large\bfseries\color{teal}}{}{0em}{}[\titlerule]
\titleformat{\subsection}{\large\bfseries\color{slate}}{}{0em}{}
\titleformat{\subsubsection}{\normalsize\bfseries\color{slate}}{}{0em}{}

%% ─── LISTINGS ────────────────────────────────────────────────────────────────
\lstset{
  backgroundcolor=\color{codebg},
  basicstyle=\footnotesize\ttfamily\color{codetext},
  breaklines=true,
  frame=single,
  rulecolor=\color{slate},
  keywordstyle=\color{teal},
  commentstyle=\color{emerald},
  stringstyle=\color{yellow!80},
  numbers=left,
  numberstyle=\tiny\color{slate},
  numbersep=8pt,
}

%% ─── TIKZ: estilos globales ──────────────────────────────────────────────────
\tikzset{
  block/.style={rectangle, draw=slate, fill=lightgray, text centered,
                rounded corners=4pt, minimum height=1cm, minimum width=3cm,
                font=\small},
  sensor/.style={block, fill=teal!20, draw=teal},
  decision/.style={block, fill=emerald!20, draw=emerald},
  actuator/.style={block, fill=orange!20, draw=orange!70},
  arrow/.style={-Latex, thick, color=slate},
  label/.style={font=\scriptsize\itshape, color=slate},
}

%%─────────────────────────────────────────────────────────────────────────────
%%  PORTADA
%%─────────────────────────────────────────────────────────────────────────────
\begin{document}

\begin{titlepage}
  \centering
  \vspace*{2cm}

  {\Huge\bfseries\color{teal} ChatDoc AI}\\[0.4cm]
  {\Large\color{slate} Agente de Consulta Documental Basado en RAG}\\[2cm]

  \begin{mdframed}[linecolor=teal, linewidth=1.5pt, innertopmargin=10pt,
                   innerbottommargin=10pt, roundcorner=6pt]
    \centering
    \textbf{Informe Técnico}\\
    Arquitectura de Agentes Inteligentes\\[0.3cm]
    {\small Identificación de Sensores, Decisiones y Actuadores}
  \end{mdframed}

  \vfill
  {\large Mayo 2026}
\end{titlepage}

%%─────────────────────────────────────────────────────────────────────────────
%% TABLA DE CONTENIDOS
%%─────────────────────────────────────────────────────────────────────────────
\tableofcontents
\newpage

%%─────────────────────────────────────────────────────────────────────────────
%% 1. RESUMEN
%%─────────────────────────────────────────────────────────────────────────────
\section{Resumen}

El presente informe describe el diseño, implementación y comportamiento del sistema \textbf{ChatDoc AI}, un agente inteligente de consulta documental que permite a los usuarios cargar documentos (PDF, DOCX, TXT) y realizar preguntas en lenguaje natural sobre su contenido. El agente aplica el paradigma \textit{Retrieval-Augmented Generation} (RAG), combinando búsqueda semántica vectorial con generación de texto mediante un modelo de lenguaje de gran escala (LLM).

La arquitectura se apoya en \textbf{LlamaIndex} como framework de orquestación, \textbf{ChromaDB} como base de datos vectorial persistente, \textbf{HuggingFace embeddings} (BAAI/bge-small-en-v1.5) para la vectorización y \textbf{Groq} con el modelo \texttt{llama-3.3-70b-versatile} como motor de respuesta. El frontend, desarrollado en React con Vite y Tailwind CSS, ofrece una interfaz de usuario moderna con soporte de tema claro/oscuro.

Se identifican y analizan los \textbf{sensores}, \textbf{módulos de decisión} y \textbf{actuadores} del agente, mostrando cómo cada componente participa en el ciclo percepción–razonamiento–acción propio de los sistemas basados en agentes.

\medskip
\textbf{Palabras clave:} RAG, LLM, agente inteligente, LlamaIndex, ChromaDB, FastAPI, React.

%%─────────────────────────────────────────────────────────────────────────────
%% 2. INTRODUCCIÓN
%%─────────────────────────────────────────────────────────────────────────────
\section{Introducción}

En los últimos años, los modelos de lenguaje de gran escala (LLMs) han demostrado capacidades notables para la comprensión y generación de texto. Sin embargo, su entrenamiento estático no contempla documentación privada o especializada propia de cada organización. El patrón \textbf{RAG} (\textit{Retrieval-Augmented Generation}) emerge como solución: combina la recuperación semántica de fragmentos relevantes de documentos con la capacidad generativa del LLM, restringiendo las respuestas al contexto cargado.

\textbf{ChatDoc AI} es un sistema RAG completo que implementa este paradigma con las siguientes características:

\begin{itemize}[leftmargin=1.5cm]
  \item Ingesta y vectorización de documentos en tiempo real (PDF, DOCX, TXT).
  \item Almacenamiento persistente de vectores con metadatos en ChromaDB.
  \item Consulta en lenguaje natural con respuestas fundamentadas en documentos.
  \item Filtrado opcional por documento específico para mayor precisión.
  \item Interfaz gráfica moderna con notificaciones de estado.
\end{itemize}

Desde la perspectiva de los agentes inteligentes, el sistema puede modelarse como un \textbf{agente reactivo informado}, donde los sensores reciben la entrada del entorno (usuario y documentos), el módulo de decisión razona sobre esa información mediante búsqueda semántica y el LLM, y los actuadores producen respuestas y modificaciones en el estado del sistema.

%%─────────────────────────────────────────────────────────────────────────────
%% 3. OBJETIVOS
%%─────────────────────────────────────────────────────────────────────────────
\section{Objetivos}

\subsection{Objetivo General}

Diseñar e implementar un agente inteligente de consulta documental que permita a los usuarios interactuar mediante lenguaje natural con el contenido de sus propios documentos, garantizando respuestas contextualmente precisas y limitadas al corpus cargado.

\subsection{Objetivos Específicos}

\begin{enumerate}[leftmargin=1.5cm]
  \item Implementar un \textit{pipeline} de ingesta documental que procese, fragmente y vectorice documentos en formato PDF, DOCX y TXT.
  \item Configurar y mantener una base de datos vectorial (ChromaDB) con persistencia en disco para almacenar los \textit{embeddings} y metadatos de los documentos.
  \item Integrar un LLM de alta capacidad (Llama 3.3-70b via Groq) para la generación de respuestas basadas exclusivamente en el contexto recuperado.
  \item Desarrollar una API REST (FastAPI) que exponga los servicios de ingesta y consulta al frontend.
  \item Construir una interfaz de usuario (React + Tailwind CSS) que permita la carga de documentos, selección de fuentes y conversación con el agente.
  \item Identificar y documentar formalmente los \textbf{sensores}, \textbf{módulos de decisión} y \textbf{actuadores} del agente para facilitar su análisis y extensión futura.
\end{enumerate}

%%─────────────────────────────────────────────────────────────────────────────
%% 4. METODOLOGÍA
%%─────────────────────────────────────────────────────────────────────────────
\section{Metodología}

El desarrollo del sistema siguió una metodología \textbf{iterativa e incremental}, organizada en las siguientes fases:

\begin{enumerate}[leftmargin=1.5cm]
  \item \textbf{Análisis de requerimientos}: Identificación de casos de uso principales (carga de documentos, consulta en lenguaje natural, filtrado por fuente) y restricciones técnicas (respuestas basadas exclusivamente en el contexto).
  \item \textbf{Diseño de arquitectura}: Selección del stack tecnológico (LlamaIndex, ChromaDB, Groq, FastAPI, React) y definición del flujo de datos entre componentes.
  \item \textbf{Implementación del backend}: Desarrollo de los servicios de ingesta (\texttt{ingesta.py}) y consulta (\texttt{chat.py}), con exposición via FastAPI.
  \item \textbf{Implementación del frontend}: Desarrollo de los componentes React (Sidebar, ChatWindow, MessageBubble, Header) y la capa de comunicación API con Axios.
  \item \textbf{Pruebas de integración}: Verificación del pipeline completo mediante el archivo \texttt{test\_chat.py} y el archivo de peticiones HTTP \texttt{requests.http}.
  \item \textbf{Análisis de comportamiento del agente}: Identificación y clasificación de sensores, decisiones y actuadores bajo el marco teórico de agentes inteligentes.
\end{enumerate}

%%─────────────────────────────────────────────────────────────────────────────
%% 5. DESARROLLO
%%─────────────────────────────────────────────────────────────────────────────
\section{Desarrollo}

\subsection{Arquitectura General del Sistema}

El sistema se estructura en dos capas principales: un \textbf{backend} Python (FastAPI) y un \textbf{frontend} JavaScript (React/Vite). Ambas capas se comunican mediante una API REST con CORS habilitado.

\begin{figure}[H]
\centering
\begin{tikzpicture}[node distance=1.4cm and 2cm]

  %% Frontend
  \node[sensor] (ui) {Interfaz de Usuario\\(React / Vite)};
  \node[block, below=0.8cm of ui] (api_client) {API Client\\(Axios)};

  %% Backend
  \node[block, right=3cm of ui] (fastapi) {FastAPI\\REST API};
  \node[sensor, below=0.8cm of fastapi] (ingesta) {Servicio Ingesta\\(ingesta.py)};
  \node[decision, right=2.5cm of fastapi] (llamaindex) {LlamaIndex\\Orchestration};

  %% Stores
  \node[actuator, below=0.8cm of llamaindex] (chromadb) {ChromaDB\\Vector Store};
  \node[decision, above=0.8cm of llamaindex] (groq) {Groq API\\Llama 3.3-70b};
  \node[block, below=0.8cm of ingesta] (hf) {HuggingFace\\bge-small-en-v1.5};

  %% Arrows
  \draw[arrow] (ui) -- (api_client);
  \draw[arrow] (api_client) -- node[above, label]{HTTP POST} (fastapi);
  \draw[arrow] (fastapi) -- (ingesta);
  \draw[arrow] (fastapi) -- (llamaindex);
  \draw[arrow] (llamaindex) -- (chromadb);
  \draw[arrow] (llamaindex) -- (groq);
  \draw[arrow] (ingesta) -- (hf);
  \draw[arrow] (hf) -- (chromadb);

\end{tikzpicture}
\caption{Diagrama de arquitectura general del sistema ChatDoc AI.}
\label{fig:arquitectura}
\end{figure}

\subsection{Identificación de Sensores, Decisiones y Actuadores}

La siguiente tabla resume de forma concisa el rol de cada componente dentro del ciclo del agente:

\begin{table}[H]
\centering
\renewcommand{\arraystretch}{1.4}
\begin{tabular}{>{\bfseries}p{2.6cm} p{3.8cm} p{5cm} p{2.8cm}}
\toprule
\textbf{Categoría} & \textbf{Componente} & \textbf{Descripción} & \textbf{Módulo} \\
\midrule
\rowcolor{teal!10}
Sensor & Entrada de texto & Pregunta del usuario en lenguaje natural & \texttt{ChatWindow.jsx} \\
\rowcolor{teal!10}
Sensor & Carga de archivo & Documento cargado (PDF/DOCX/TXT) via drag-and-drop & \texttt{Sidebar.jsx} \\
\rowcolor{teal!10}
Sensor & Selección de fuente & Archivos activos para filtrar la consulta & \texttt{Sidebar.jsx} \\
\rowcolor{teal!10}
Sensor & Estado HTTP & Respuesta del servidor (éxito/error de conexión) & \texttt{api.js} \\
\midrule
\rowcolor{emerald!10}
Decisión & Vectorización & Conversión de texto a embeddings semánticos & \texttt{ingesta.py} \\
\rowcolor{emerald!10}
Decisión & Indexación & Almacenamiento vectorial con metadatos en ChromaDB & \texttt{ingesta.py} \\
\rowcolor{emerald!10}
Decisión & Recuperación semántica & Búsqueda por similitud del vector de la pregunta & \texttt{chat.py} \\
\rowcolor{emerald!10}
Decisión & Filtrado por fuente & Aplicación de MetadataFilters por nombre de archivo & \texttt{chat.py} \\
\rowcolor{emerald!10}
Decisión & Generación RAG & Síntesis de respuesta con LLM + contexto recuperado & \texttt{chat.py} \\
\rowcolor{emerald!10}
Decisión & Prompt restrictivo & Política de respuesta: solo contexto, sin inventar & \texttt{chat.py} \\
\midrule
\rowcolor{orange!10}
Actuador & Respuesta en chat & Texto de respuesta mostrado en la interfaz & \texttt{ChatWindow.jsx} \\
\rowcolor{orange!10}
Actuador & Notificación toast & Alerta visual de éxito o error de ingesta & \texttt{Sidebar.jsx} \\
\rowcolor{orange!10}
Actuador & Escritura en ChromaDB & Persistencia de vectores y documentos indexados & \texttt{ingesta.py} \\
\rowcolor{orange!10}
Actuador & Indicador de carga & Animación visual mientras el agente procesa & \texttt{ChatWindow.jsx} \\
\bottomrule
\end{tabular}
\caption{Clasificación de sensores, decisiones y actuadores del agente ChatDoc AI.}
\label{tab:sda}
\end{table}

\subsection{Sensores del Agente}

Los sensores son los mecanismos mediante los cuales el agente \textbf{percibe el entorno}. En ChatDoc AI se identifican cuatro sensores principales:

\subsubsection{S1 -- Entrada de Texto (Pregunta del Usuario)}

El componente \texttt{ChatWindow.jsx} contiene un área de texto (\texttt{<textarea>}) que captura la pregunta del usuario. La acción se dispara al presionar \texttt{Enter} (sin Shift) o el botón de envío. El texto es validado (no vacío, no en estado de carga) antes de enviarse.

\subsubsection{S2 -- Carga de Documento}

El componente \texttt{Sidebar.jsx} implementa una zona de \textit{drag-and-drop} mediante la librería \texttt{react-dropzone}. Acepta los tipos MIME: \texttt{application/pdf}, \texttt{text/plain}, \texttt{application/vnd.openxmlformats-officedocument.wordprocessingml.document} y \texttt{application/msword}. Cada archivo se envía al endpoint \texttt{POST /ingestar/} del backend.

\subsubsection{S3 -- Selección de Fuente Activa}

El usuario puede marcar/desmarcar documentos en la lista de fuentes del sidebar. Esta selección determina si la consulta se aplica a \textbf{todos los documentos} o solo a los \textbf{archivos seleccionados}, afectando directamente al filtrado semántico.

\subsubsection{S4 -- Estado de la Respuesta HTTP}

La capa de API (\texttt{api.js}) utiliza Axios para comunicarse con el backend. El resultado de las llamadas (éxito o error HTTP) es interpretado por el frontend para actualizar el estado de la interfaz y mostrar notificaciones adecuadas.

\subsection{Módulos de Decisión del Agente}

Los módulos de decisión implementan el \textbf{razonamiento} del agente: cómo procesa la percepción para determinar una acción.

\subsubsection{D1 -- Vectorización con HuggingFace Embeddings}

Durante la ingesta, cada fragmento del documento es convertido a un vector numérico de alta dimensión mediante el modelo \texttt{BAAI/bge-small-en-v1.5}. Este modelo de código abierto, ejecutado localmente, captura el significado semántico del texto.

\begin{lstlisting}[language=Python, caption={Configuración del modelo de embeddings.}]
Settings.embed_model = HuggingFaceEmbedding(
    model_name="BAAI/bge-small-en-v1.5"
)
\end{lstlisting}

\subsubsection{D2 -- Indexación en ChromaDB}

Los vectores generados son almacenados en una colección persistente de ChromaDB (\texttt{documentos\_usuario}), junto con los metadatos del documento (nombre de archivo, fragmento de texto). Este índice persiste en disco en \texttt{./data/chroma\_db}.

\subsubsection{D3 -- Recuperación Semántica (Similarity Search)}

Al recibir una pregunta, el agente vectoriza la consulta con el mismo modelo de embeddings y realiza una \textbf{búsqueda por similitud del coseno} en ChromaDB, recuperando los fragmentos más relevantes del corpus.

\subsubsection{D4 -- Filtrado por Metadatos}

Si el usuario ha seleccionado uno o más documentos específicos, se aplican \texttt{MetadataFilters} con condición \texttt{OR} sobre el campo \texttt{file\_name}, restringiendo la búsqueda semántica solo a los archivos indicados.

\begin{lstlisting}[language=Python, caption={Aplicación de filtros por archivo en la consulta.}]
filtros = [
    MetadataFilter(key="file_name", value=archivo,
                   operator=FilterOperator.EQ)
    for archivo in archivos
]
filters = MetadataFilters(
    filters=filtros, condition=FilterCondition.OR
)
\end{lstlisting}

\subsubsection{D5 -- Generación de Respuesta (RAG con Groq/Llama 3.3)}

Los fragmentos recuperados se ensamblan como contexto y se envían junto con la pregunta al modelo \texttt{llama-3.3-70b-versatile} via la API de Groq. El LLM genera la respuesta final.

\subsubsection{D6 -- Prompt Restrictivo (Control de Alucinaciones)}

Un \textit{system prompt} estricto instruye al modelo a responder \textbf{únicamente} con información del contexto proporcionado. Si la respuesta no está disponible en los documentos, el agente responde: \textit{``Lo siento, esa información no se encuentra en los documentos proporcionados.''}

\subsection{Actuadores del Agente}

Los actuadores son los mecanismos por los cuales el agente \textbf{modifica el entorno}.

\subsubsection{A1 -- Respuesta en Lenguaje Natural}

La respuesta generada por el LLM es recibida por el frontend y renderizada en el componente \texttt{MessageBubble.jsx} con soporte Markdown, diferenciando visualmente mensajes del usuario (derecha) y del asistente (izquierda).

\subsubsection{A2 -- Notificaciones Toast}

El sistema emite notificaciones visuales no intrusivas mediante la librería \texttt{react-hot-toast}: notificación de éxito al completar la ingesta de un documento, y notificación de error en caso de fallo.

\subsubsection{A3 -- Escritura en Base de Datos Vectorial}

La acción más relevante desde el punto de vista del estado interno del agente: la función \texttt{procesar\_y\_almacenar()} escribe los vectores e índices en ChromaDB de forma persistente, ampliando el conocimiento del agente con cada nuevo documento cargado.

\subsubsection{A4 -- Indicador Visual de Procesamiento}

Durante la espera de la respuesta, el componente \texttt{ChatWindow.jsx} muestra una animación de tres puntos pulsantes que informa al usuario que el agente está procesando su consulta.

\subsection{Ciclo del Agente: Diagrama de Flujo}

\begin{figure}[H]
\centering
\begin{tikzpicture}[node distance=1.3cm, every node/.style={font=\small}]

  \node[sensor, minimum width=4.5cm] (s1) {[S1] Usuario escribe pregunta};
  \node[sensor, minimum width=4.5cm, below=0.7cm of s1] (s3) {[S3] Archivos seleccionados};
  \node[decision, minimum width=4.5cm, below=0.7cm of s3] (d3) {[D3] Vectorizar pregunta\\+ Búsqueda semántica};
  \node[decision, minimum width=4.5cm, below=0.7cm of d3] (d4) {[D4] ¿Filtro por archivo?};
  \node[decision, minimum width=4.5cm, below=0.7cm of d4] (d5) {[D5] RAG: Contexto + LLM};
  \node[actuator, minimum width=4.5cm, below=0.7cm of d5] (a1) {[A1] Mostrar respuesta};

  \node[sensor, minimum width=4.5cm, right=3.5cm of s1] (s2) {[S2] Usuario carga documento};
  \node[decision, minimum width=4.5cm, below=0.7cm of s2] (d1) {[D1] Vectorizar fragmentos\\(HuggingFace)};
  \node[decision, minimum width=4.5cm, below=0.7cm of d1] (d2) {[D2] Indexar en ChromaDB};
  \node[actuator, minimum width=4.5cm, below=0.7cm of d2] (a3) {[A3] Persistir vectores};
  \node[actuator, minimum width=4.5cm, below=0.7cm of a3] (a2) {[A2] Notificación toast};

  \draw[arrow] (s1) -- (s3);
  \draw[arrow] (s3) -- (d3);
  \draw[arrow] (d3) -- (d4);
  \draw[arrow] (d4) -- (d5);
  \draw[arrow] (d5) -- (a1);

  \draw[arrow] (s2) -- (d1);
  \draw[arrow] (d1) -- (d2);
  \draw[arrow] (d2) -- (a3);
  \draw[arrow] (a3) -- (a2);

  %% Leyenda
  \node[sensor, minimum width=1.5cm, right=0.5cm of a3, yshift=1.5cm] (ls) {Sensor};
  \node[decision, minimum width=1.5cm, below=0.4cm of ls] (ld) {Decisión};
  \node[actuator, minimum width=1.5cm, below=0.4cm of ld] (la) {Actuador};

\end{tikzpicture}
\caption{Flujo de percepción-razonamiento-acción del agente ChatDoc AI.}
\label{fig:flujo}
\end{figure}

\subsection{Stack Tecnológico}

\begin{table}[H]
\centering
\renewcommand{\arraystretch}{1.3}
\begin{tabular}{llll}
\toprule
\textbf{Capa} & \textbf{Tecnología} & \textbf{Versión/Variante} & \textbf{Rol} \\
\midrule
Frontend & React & 18+ & UI / Estado \\
Frontend & Vite & 5+ & Bundler / Dev server \\
Frontend & Tailwind CSS & 3 & Estilos \\
Frontend & Framer Motion & -- & Animaciones \\
Frontend & Axios & -- & Cliente HTTP \\
Backend & Python & 3.10+ & Lenguaje principal \\
Backend & FastAPI & -- & API REST \\
Backend & LlamaIndex & -- & Orquestación RAG \\
Backend & ChromaDB & -- & Base de datos vectorial \\
Backend & HuggingFace & bge-small-en-v1.5 & Embeddings \\
Backend & Groq API & llama-3.3-70b-versatile & LLM \\
Backend & PyMuPDF & -- & Extracción PDF \\
\bottomrule
\end{tabular}
\caption{Stack tecnológico completo de ChatDoc AI.}
\label{tab:stack}
\end{table}

%%─────────────────────────────────────────────────────────────────────────────
%% 6. RESULTADOS
%%─────────────────────────────────────────────────────────────────────────────
\section{Resultados}

\subsection{Comportamiento del Agente}

El agente ChatDoc AI exhibe el siguiente comportamiento observable durante su operación:

\begin{enumerate}[leftmargin=1.5cm]
  \item \textbf{Fase de ingesta}: Al cargar un documento, el agente lo fragmenta automáticamente, genera \textit{embeddings} para cada fragmento y los persiste en ChromaDB. El proceso es transparente para el usuario, quien recibe una notificación visual de éxito o fallo. El backend confirma el número de documentos procesados en la respuesta JSON.

  \item \textbf{Fase de consulta}: Al formular una pregunta, el agente:
    \begin{enumerate}[label=\alph*.]
      \item Vectoriza la pregunta con el mismo modelo de embeddings.
      \item Recupera los fragmentos más similares del índice (con o sin filtro de archivo).
      \item Construye un prompt enriquecido con el contexto recuperado.
      \item Genera la respuesta con el LLM y la devuelve al frontend.
    \end{enumerate}

  \item \textbf{Comportamiento ante preguntas fuera de contexto}: Cuando la pregunta no puede responderse con los documentos cargados, el agente responde con el mensaje predefinido: \textit{``Lo siento, esa información no se encuentra en los documentos proporcionados.''}, evitando alucinaciones.

  \item \textbf{Filtrado por fuente}: Cuando el usuario selecciona archivos específicos en el sidebar, el agente restringe la búsqueda semántica únicamente a esos documentos, devolviendo respuestas más precisas y acotadas.
\end{enumerate}

\subsection{Identificación Final: Tabla Consolidada SDA}

\begin{longtable}{p{1.5cm} p{1.5cm} p{3.5cm} p{5.5cm}}
\toprule
\textbf{ID} & \textbf{Tipo} & \textbf{Nombre} & \textbf{Función} \\
\midrule
\endfirsthead
\toprule
\textbf{ID} & \textbf{Tipo} & \textbf{Nombre} & \textbf{Función} \\
\midrule
\endhead
S1 & Sensor & Entrada de texto & Captura la pregunta del usuario vía teclado o botón \\
S2 & Sensor & Carga de archivo & Recibe documentos por drag-and-drop o selector \\
S3 & Sensor & Selección de fuente & Indica qué archivos usar como contexto de búsqueda \\
S4 & Sensor & Estado HTTP & Detecta éxito o fallo en la comunicación con el backend \\
\midrule
D1 & Decisión & Vectorización & Transforma texto en vectores semánticos (HuggingFace) \\
D2 & Decisión & Indexación & Almacena vectores y metadatos en ChromaDB \\
D3 & Decisión & Recuperación semántica & Busca fragmentos relevantes por similitud de coseno \\
D4 & Decisión & Filtrado por metadatos & Restringe la búsqueda por nombre de archivo \\
D5 & Decisión & Generación RAG & Produce respuesta con Groq/Llama 3.3 y contexto \\
D6 & Decisión & Prompt restrictivo & Previene alucinaciones; responde solo del corpus \\
\midrule
A1 & Actuador & Respuesta en chat & Muestra la respuesta generada en la interfaz \\
A2 & Actuador & Notificación toast & Informa sobre el resultado de la ingesta \\
A3 & Actuador & Escritura en ChromaDB & Persiste los vectores del nuevo documento \\
A4 & Actuador & Indicador de carga & Muestra animación mientras el agente procesa \\
\bottomrule
\caption{Tabla consolidada de Sensores, Decisiones y Actuadores del agente ChatDoc AI.}
\label{tab:sda_final}
\end{longtable}

%%─────────────────────────────────────────────────────────────────────────────
%% 7. CONCLUSIONES
%%─────────────────────────────────────────────────────────────────────────────
\section{Conclusiones}

\begin{enumerate}[leftmargin=1.5cm]
  \item ChatDoc AI implementa exitosamente el paradigma RAG, logrando un agente capaz de responder preguntas en lenguaje natural sobre documentos propios del usuario, sin recurrir a conocimiento externo al corpus proporcionado.

  \item El análisis de sensores, decisiones y actuadores confirma que el sistema opera como un \textbf{agente reactivo informado}: sus sensores perciben tanto el input del usuario (pregunta, archivo) como el estado del sistema (selección activa, estado HTTP); sus módulos de decisión aplican razonamiento semántico (embeddings, recuperación, filtrado, generación) y sus actuadores modifican el entorno de forma útil (respuesta, persistencia, notificación).

  \item La combinación de embeddings locales (HuggingFace) con un LLM en la nube (Groq) representa un balance efectivo entre costo computacional y calidad de respuesta.

  \item El uso de ChromaDB con persistencia garantiza que el conocimiento indexado no se pierde entre sesiones, permitiendo un crecimiento progresivo de la base de conocimiento del agente.

  \item El prompt restrictivo implementado como D6 es crítico para la confiabilidad del sistema: elimina el riesgo de que el agente genere respuestas falsas basadas en su conocimiento paramétrico.
\end{enumerate}

%%─────────────────────────────────────────────────────────────────────────────
%% 8. RECOMENDACIONES
%%─────────────────────────────────────────────────────────────────────────────
\section{Recomendaciones}

\begin{enumerate}[leftmargin=1.5cm]
  \item \textbf{Autenticación y control de acceso}: Implementar un sistema de autenticación (JWT, OAuth2) para aislar los documentos por usuario, especialmente si el sistema se despliega en un entorno multiusuario.

  \item \textbf{Gestión de colecciones por usuario}: Crear colecciones ChromaDB independientes por usuario o sesión, evitando que las preguntas de un usuario accedan a documentos de otro.

  \item \textbf{Eliminación de documentos}: Agregar la funcionalidad de eliminar documentos del índice, incluyendo la remoción de sus vectores en ChromaDB (actualmente no existe un endpoint para esta operación).

  \item \textbf{Chunking configurable}: Exponer parámetros de fragmentación (tamaño de chunk, overlap) como opciones de configuración, ya que tienen impacto directo en la calidad de la recuperación semántica.

  \item \textbf{Memoria conversacional}: Integrar memoria de contexto conversacional (historial de turnos) en el módulo de decisión D5 para que el agente pueda responder preguntas de seguimiento sin perder el hilo de la conversación.

  \item \textbf{Soporte de imágenes y tablas en PDF}: Mejorar el extractor de PDF (actualmente PyMuPDF) con capacidades multimodales para indexar contenido visual dentro de los documentos.

  \item \textbf{Evaluación de calidad RAG}: Incorporar métricas de evaluación (fidelidad, relevancia de recuperación) usando frameworks como \texttt{RAGAS} para monitorear la calidad del agente de forma continua.

  \item \textbf{Despliegue con Docker}: Containerizar el backend y el frontend con Docker Compose para facilitar el despliegue reproducible en cualquier entorno, incluyendo la correcta persistencia del volumen de ChromaDB.
\end{enumerate}

%%─────────────────────────────────────────────────────────────────────────────
%% FIN DEL DOCUMENTO
%%─────────────────────────────────────────────────────────────────────────────
\end{document}
